% Chapter 3.5: Cardiac Electrophysiology and Arrhythmia Dynamics
% Part III of the book.

\chapter{Cardiac Electrophysiology and Arrhythmia Dynamics}
\label{ch:part3-cardiac-electrophysiology-and-arrhythmia-d}
\label{sec:cardiac}
% -----------------------------------------------------------------------------

Most sudden cardiac deaths are not caused by a heart that stops, but by one that
quivers. Ventricular fibrillation\index{ventricular fibrillation} is an electrical
turbulence: instead of a single wavefront recruited by the pacemaker and swept
through both ventricles in about as long as it takes the chamber to squeeze, the
tissue is occupied by several reentrant spiral waves of transmembrane voltage
that break, meander and annihilate one another. Pumping fails because no
sufficient fraction of the myocardium is in systole at the same moment. The
disease is electrical, so the interventions that matter are electrical and
spatial---where to ablate, which tissue parameters explain a recorded
rhythm---and each begins with the question \emph{what does this tissue do next,
given this starting state}. That is an operator question.

\index{cardiac electrophysiology}The governing system is a stiff
reaction--diffusion equation on an anisotropic domain with the shape of a
ventricle. Its solutions contain two features that a numerical method must pay
for: a depolarizing front that rises almost discontinuously, and a recovery
process that takes hundreds of milliseconds behind
it. Resolving the first fixes the mesh and the time step; simulating long enough
to see a rotor drift fixes the number of steps. The per-solve cost of
\chapref{ch:part3-the-multi-query-bottleneck-in-computationa} is therefore not
incidental: workflows needing many solves at once---parameter inference from
body-surface potentials, lesion-set search, ensembles under perturbed
stimulation---inherit it in full.

What Parts~I and~II bring to that cost is a small set of precise instruments:
the composite Physics-Informed Neural Operator (PINO) objective \eqref{eq:pino_loss}
of \chapref{ch:part1-physics-informed-operators-and-sobolev-reg},
the exact zero-shot identity of
\chapref{ch:part1-discretization-invariance-and-continuum-li} with its
band-limitedness hypothesis, the recursive-rollout accounting of
\eqref{eq:ch09-gronwall}, and the pushforward training recipe
\eqref{eq:ch26-pushforward} of \chapref{ch:part2-training-regimes-optimization-and-failure-}.
Cardiac tissue is an awkward place to test them. The front is the least
band-limited object in the field, the monodomain operator is not coercive in any
useful sense so the a priori estimate \eqref{eq:ch07-apriori} does not apply,
there is no conservation law to check the answer against, and the labels are
simulator output rather than measurements. This chapter states the model and its
unit conventions, formulates the time-$\Delta t$ map a surrogate should learn, and
writes out what the PINO template \cite{li2021physics} becomes when it is applied
to that map. The route runs from \eqref{eq:monodomain} to the figure of a
rotor, and from there to a table of what is documented, what is asserted only by
this Part's working record, and what is not documented at all --- a column that, for the
surrogate side of this ledger, is now most of it.

% -----------------------------------------------------------------------------
\section{The Monodomain Reaction-Diffusion System}
\label{sec:ch35-monodomain}
\index{monodomain equation}
% -----------------------------------------------------------------------------

Electrical propagation across cardiac muscle is modeled by the monodomain
equation:
\begin{equation}
  \label{eq:monodomain}
  \chi \left( C_m \frac{\partial V_m}{\partial t} + I_{\text{ion}}(V_m, \mathbf{w}) \right) = \nabla \cdot (\mathbf{D}(x) \nabla V_m) + I_{\text{stim}}(x, t),
\end{equation}
where $V_m(x, t)$ is transmembrane voltage\index{$V_m$@$V_m$, transmembrane voltage}, $C_m$ is membrane capacitance\index{$C_m$@$C_m$, membrane capacitance}, $\chi$ is surface-to-volume ratio\index{$\chi$@$\chi$, surface-to-volume ratio}, and $\mathbf{D}(x)$ is the diffusion tensor reflecting electrical anisotropy along myocardial fibers:
\begin{equation}
  \label{eq:ch35-dtensor}
  \mathbf{D}(x) = \sigma_t \mathbf{I} + (\sigma_l - \sigma_t) \mathbf{f}_0(x) \otimes \mathbf{f}_0(x),
\end{equation}
where $\mathbf{f}_0(x)$ is the local fiber orientation vector\index{fiber orientation field}\index{cardiac anisotropy}. The ion channel current $I_{\text{ion}}$ is parameterized by reaction kinetics such as the two-variable Aliev-Panfilov model\index{Aliev-Panfilov model}~\cite{aliev1996simple}:
\begin{align}
  \label{eq:ch35-ion}
  I_{\text{ion}}(V_m, w) &= k V_m (V_m - a)(V_m - 1) + V_m w, \\
  \label{eq:ch35-recovery}
  \frac{\partial w}{\partial t} &= \big( \epsilon_0 + \frac{\mu_1 w}{V_m + \mu_2} \big) \big( -w - k V_m (V_m - b - 1) \big).
\end{align}
Boundary conditions require zero electrical flux through the endocardium and epicardium:
\begin{equation}
  \label{eq:ch35-no-flux}
  \hat{\mathbf{n}} \cdot (\mathbf{D} \nabla V_m) = 0 \quad \text{on } \partial\Omega.
\end{equation}

Each symbol in \eqref{eq:monodomain} carries a modeling decision, and the
decisions matter more than the notation. The left-hand side is a current balance
per unit membrane area: $C_m\,\partial_t V_m$ is the capacitive current charging
that area, and $I_{\text{ion}}$ the ionic current carried through membrane
channels; $\chi$ converts per-area quantities into per-volume quantities for a
syncytium of cells, and the surface-to-volume ratio is why a model of myocytes
and a model of tissue are the same equation with different constants. On the
right, $\nabla\cdot(\mathbf{D}\nabla V_m)$ is the axial current flowing through
the gap junctions between cells, and $I_{\text{stim}}$ is the applied stimulus,
either a pacing electrode or the shock of a defibrillator. The reduction to a
single voltage field assumes the extracellular space is grounded and that its
anisotropy ratio matches the intracellular one; where the extracellular
potential itself is wanted---body-surface electrocardiogram (ECG) synthesis,
field stimulation---that assumption fails and the two-domain bidomain system
must be solved instead.\index{bidomain model}\index{electrocardiogram (ECG)}
The monodomain equation is thus a deliberate narrowing: cheaper, well posed
enough for wave propagation studies, and silent about the extracellular
recordings that are the clinician's only window on the organ.

\begin{remark}[Normalization and units]
\label{rem:ch35-units}
Equations \eqref{eq:monodomain} and \eqref{eq:ch35-ion} do not share a unit
system as written, and the discrepancy is instructive rather than cosmetic. The
cubic factor $V_m(V_m-a)(V_m-1)$ in \eqref{eq:ch35-ion} has zeros at $0$ and
$1$, which is the standard nondimensionalization in which $V_m$ is scaled to
the interval between the resting potential and the overshoot of the upstroke; in
that scaling $V_m$ is dimensionless, $a$ and $b$ are dimensionless shape
parameters, and $k,\epsilon_0,\mu_1,\mu_2$ are dimensionless kinetic constants.
A dimensional $I_{\text{ion}}$ in $\mu\text{A}/\text{cm}^2$ added to
$C_m\,\partial_t V_m$ is a different object. The physiological reading of
\eqref{eq:monodomain} is
that $V_m$ is a membrane potential in millivolts, negative at rest and positive
at the top of the upstroke; the numerical reading, and the one the reaction
model \eqref{eq:ch35-ion}--\eqref{eq:ch35-recovery} presumes, is a
dimensionless excitation variable. The bridge has to be declared rather than
inherited. Taking rest at $-80\ \mathrm{mV}$ and the overshoot at $+20\ \mathrm{mV}$,
the scaling that places the zeros of the cubic at the two endpoints is
$V\mapsto(V[\mathrm{mV}]+80)/100$, and time must be divided by a chosen
characteristic excitation time of the tissue. Those two divisors are precisely
the voltage and time scales that
have to be restored before a dimensional $I_{\text{ion}}$ can be added to
$C_m\,\partial_t V_m$. Any residual penalty built
from \eqref{eq:monodomain} inherits this ambiguity, and because the weights of
\eqref{eq:pino_loss} are dimensional---\chapref{ch:part1-physics-informed-operators-and-sobolev-reg}
says so explicitly---the choice of scaling changes which $\lambda_{\text{pde}}$
balances the terms. This is the place in Part~III where that observation stops
being abstract.
\end{remark}

An Aliev-Panfilov current \eqref{eq:ch35-ion} is not a description of ion
channels. It is the minimal excitable structure that reproduces the phenomena a
wave-propagation study needs: a threshold separating a stimulus that dies out
from one that triggers a full action potential, a refractory tail governed by
the slow recovery variable $w$ in \eqref{eq:ch35-recovery}, a propagation speed
set by the ratio of the diffusion coefficient to the excitability, and the
restitution relation between diastolic interval and action potential duration
(APD)\index{action potential duration (APD)} that organizes alternans and
breakup. What it discards is ionic specificity: no separate sodium, calcium and
potassium currents, therefore no mechanism by which a drug or a mutation
prolongs the action potential. Fully ionic models of a single human ventricular
myocyte, carrying scores of state variables, exist for exactly that purpose,
and a surrogate trained on \eqref{eq:ch35-ion} reaches them only by being
retrained on them. The number of state variables is the first quantity to read
off any cardiac operator-learning claim, because it is the number of output
channels the network must carry and the number of equations the residual term
must satisfy; for the model \eqref{eq:ch35-ion}--\eqref{eq:ch35-recovery} it is
two, the voltage and a recovery field.

\index{stiff reaction--diffusion system}Cost is where the structure of \eqref{eq:monodomain} becomes a budget. Write the
diffusivity as $\mathcal{D}(x)=\mathbf{D}(x)/(\chi C_m)$ and let
$D_{\max}=\max_x \rho_{\max}(\mathcal{D}(x))$ be the largest eigenvalue of that
tensor anywhere in the tissue. A forward-Euler treatment of the diffusion
operator on a uniform grid of spacing $h$ in $d$ spatial dimensions is stable
only while
\begin{equation}
  \label{eq:ch35-cfl}
  \Delta t \;\le\; \frac{h^{2}}{2\,d\,D_{\max}},
\end{equation}
the diffusion restriction of the more familiar Courant--Friedrichs--Lewy
condition, and the truncation error of any fixed-order difference scheme in
space and time behaves the same way, shrinking algebraically in $h$ and $\Delta
t$~\cite{leveque2002finite,quarteroni2008numerical}. Implicit stepping buys
stability at $h^2$-scale steps, but the reaction term \eqref{eq:ch35-ion} is
stiff on the upstroke, so the practical step is still set by accuracy through
the front rather than by the linear solve. Integrating to a simulated horizon
$T$ therefore costs, in the explicit case,
\begin{equation}
  \label{eq:ch35-cost}
  W(h,T) \;=\; \frac{T}{\Delta t}\,\mathcal{O}(N_h)
  \;=\; \mathcal{O}\!\left(T\,h^{-(d+2)}\right),
  \qquad N_h \propto h^{-d},
\end{equation}
where $N_h$ is the number of unknowns. Two consequences organize the rest of
this chapter. Refinement is punitive: halving $h$ in three dimensions multiplies
the work by roughly $32$, and the exponent in \eqref{eq:ch35-cost} is the reason
the inherited mesh width and time step, $h\sim 200\ \mu\text{m}$ and
$\Delta t\sim 10\ \mu\text{s}$, should be read as a coupled pair rather than as
two independent choices. And the front sets the mesh: nodes must be spaced a
small fraction of a millimeter wherever a wave is going to travel, which over an
anatomical ventricle is what puts a monodomain solve in the
$N\sim10^{7}$--$10^{9}$ regime of
\chapref{ch:part3-the-multi-query-bottleneck-in-computationa}. Those two figures
are order-of-magnitude values carried forward into this Part; they are customary
for monodomain work discretized by the finite element method (FEM), of the kind
the openCARP simulator performs, but this book cites no primary measurement of
them, and Table~\ref{tab:ch35cost} says so. The regime that count describes is
also the one no operator surrogate reaches anywhere in this book's bibliography.
The distinction is dimensional: on a flat two-dimensional sheet of linear size
$L$ meshed at $h$ the unknowns scale as $(L/h)^2$ rather than $(L/h)^3$, and with
$V_m$ and $w$ the state carries two fields rather than the scores of a fully
ionic model. A sheet of side $10\ \text{cm}$ discretized at the same
$h\sim200\ \mu\text{m}$ quoted above carries about $2.5\times10^{5}$ nodes, two to four orders
below the whole-organ count above, and the PINO formulation
\cite{li2021physics} of Section~\ref{sec:ch35-map} is posed at that smaller size.
The gap is not an
objection to the template, which tests an architecture rather than a pipeline; it
is the reason the accounting in Section~\ref{sec:ch35-claims} has to be done row
by row.
\index{finite element method (FEM)!cardiac monodomain discretization}
\index{Courant--Friedrichs--Lewy condition!diffusion restriction}

% -----------------------------------------------------------------------------
\section{\texorpdfstring{Learning the Time-$\Delta t$ Map of an Excitable Medium}{Learning the Time-Delta t Map of an Excitable Medium}}
\label{sec:ch35-map}
\index{time-$\Delta t$ map!cardiac}
% -----------------------------------------------------------------------------

The object a surrogate should provide is not the solution of
\eqref{eq:monodomain} as a function of space and time but one step of its
dynamics, applied repeatedly:
\begin{equation}
  \label{eq:ch35-timemap}
  \mathcal{G}_{\Delta t}:\;
  \big(V_m(\cdot,t_k),\,\mathbf{w}(\cdot,t_k);\,\mathbf{f}_0,\,\sigma_l,\,\sigma_t,\,I_{\text{stim}}(\cdot,t_k)\big)
  \;\longmapsto\;
  \big(V_m(\cdot,t_k+\Delta t),\,\mathbf{w}(\cdot,t_k+\Delta t)\big).
\end{equation}
Equation \eqref{eq:ch35-timemap} is the same recursion as \eqref{eq:ch26-rollout},
with two domain-specific
differences. First, the state is a pair of fields---excitation and
recovery---so the output channel count is at least two, and a surrogate that
predicts only $V_m$ cannot be advanced by itself. Second, the fiber orientation
field $\mathbf{f}_0$ and the conductivities in \eqref{eq:ch35-dtensor} are
\emph{inputs} of the map, not constants of the architecture: making them channels
of the input function is what allows one trained network to serve many patients
and many anisotropy assumptions, which is precisely the operator-learning
argument of \chapref{ch:part1-partial-differential-equations-as-solution}
applied to tissue instead of to a torus.

\index{Physics-Informed Neural Operator (PINO)!cardiac electrophysiology}\index{Fourier neural operator (FNO)!cardiac surrogate}The
architecture that map calls for is a \emph{Fourier Neural
Operator} (FNO) backbone~\cite{li2021fourier} carrying the step, with the cell
model \eqref{eq:ch35-ion}--\eqref{eq:ch35-recovery} imposed as a soft
constraint rather than solved: a PINO in the sense of \cite{li2021physics}. The
formulation is narrower than \eqref{eq:ch35-timemap} as soon as it is written
down, and the difference is worth recording precisely. The benchmark problems
PINO is demonstrated on \cite{li2021physics} have scalar diffusion coefficients
on fixed grids, so \eqref{eq:ch35-dtensor} collapses in that family and nothing
in it tests whether a fiber field can be carried as an argument of the operator.
The patient-varying reading of \eqref{eq:ch35-timemap}, in which $\mathbf{f}_0$
and the two conductivities are arguments of the operator, is this book's
formulation, and no study cited here demonstrates
it.\index{cardiac anisotropy!as an operator input}
Read against
\eqref{eq:pino_loss}, the design is exactly the template of
\chapref{ch:part1-physics-informed-operators-and-sobolev-reg}: a data term over
reference trajectories, a residual term measuring violation of
\eqref{eq:monodomain}, and a boundary term for the no-flux condition
\eqref{eq:ch35-no-flux}. Writing
the strong monodomain residual as
\begin{equation}
  \label{eq:ch35-residual}
  r_{\mathcal{N}}(V_m,\mathbf{w})
  \;=\; \chi\Big(C_m\,\partial_t V_m + I_{\text{ion}}(V_m,\mathbf{w})\Big)
  - \nabla\cdot\!\big(\mathbf{D}\nabla V_m\big) - I_{\text{stim}},
\end{equation}
the composite objective has the shape
\begin{equation}
  \label{eq:ch35-loss}
  \mathcal{L}(\theta)
  = \lambda_{\text{data}}\,\mathcal{L}_{\text{data}}(\theta)
  + \lambda_{\text{pde}}\,\big\|r_{\mathcal{N}}\big(\mathcal{G}_\theta(a)\big)\big\|_{L^2(\Omega_T)}^2
  + \lambda_{\text{bc}}\,\big\|\hat{\mathbf{n}}\cdot\mathbf{D}\nabla \mathcal{G}_\theta(a)\big\|_{L^2(\partial\Omega\times(0,T))}^2 .
\end{equation}
The first term is \eqref{eq:ch07-data-loss}, optionally in the pushforward form
\eqref{eq:ch26-pushforward}; the second is the soft physics penalty that
\eqref{eq:ch35-recovery} and \eqref{eq:ch35-ion} make unusually informative,
because the reaction kinetics are known exactly even when the tissue parameters
are not; the third penalizes violation of \eqref{eq:ch35-no-flux}, and it is no
afterthought. Part of it is forced rather than chosen: the homogeneous Neumann
condition has to be imposed on the borders of the predicted grid, and the field
must be padded there so that the periodicity assumed inside the FNO layers is not
violated --- the same device
\chapref{ch:part3-subsurface-multiphase-flow-and-carbon-stor} meets at a
reservoir wall,
where padding is a workaround for the torus hypothesis of \eqref{eq:fno_layer}
and not a discharge of it. What padding costs is what a spectral kernel gives up
on a body that is not periodic, and it is why
\chapref{ch:part3-boundary-condition-enforcement-neural-oper}
returns to this example. What remains free is the relative size of
$\lambda_{\text{pde}}$ and $\lambda_{\text{bc}}$, and
\chapref{ch:part1-physics-informed-operators-and-sobolev-reg} has already said
why that freedom is not cosmetic: the weights are dimensional.

Why penalize a residual at all here, when a solver can generate labels? Because
the residual supplies constraints at the mesh points where supervision runs out.
\chapref{ch:part1-physics-informed-operators-and-sobolev-reg} makes the general
point---an input together with the equation supplies a target---and in this
domain the specific payoff is the sub-grid field. A prediction evaluated on a
grid finer than the training grid is unconstrained by data at the added nodes;
$r_{\mathcal{N}}$ evaluated there is not, since \eqref{eq:ch35-residual} contains
$\partial_t V_m$ and $\nabla\cdot(\mathbf{D}\nabla V_m)$, both computable at any
resolution the storage supports. That is a mechanism, not a guarantee, and it
belongs next to the qualification of Section~\ref{sec:ch35-machinery}.

% -----------------------------------------------------------------------------
\section{Residual Machinery, Zero-Shot Transfer, and the Rotor Geometry}
\label{sec:ch35-machinery}
\index{spectral differentiation!at a moving front}
% -----------------------------------------------------------------------------

The machinery that evaluates
\eqref{eq:ch35-residual} cheaply is \eqref{eq:spectral_diff}: derivatives come
from multiplying Fourier coefficients by $(2\pi i k_j)^m$, at two traversals of
the transform pipeline per derivative. But the qualification Part~I attaches to
that identity is the one that bites here. The symbol has magnitude
$(2\pi|k|)^m$, so content near the Nyquist wavenumber is amplified rather than
damped---and the depolarizing front is precisely the component of $V_m$ whose
spectrum reaches that far. A Phase 0 upstroke is a moving internal layer, not a
band-limited function; the residual built from spectral derivatives of a
network's own interpolant therefore measures the sharpest thing the network has
not learned smoothly. Two further difficulties travel with it: the aliasing of
the nonlinear products $V_m^3$ and $V_m w$ in \eqref{eq:ch35-ion} returns the
dealiasing requirement discussed in Part~I, and on a ventricle-shaped domain
with the homogeneous Neumann condition \eqref{eq:ch35-no-flux} there is no
periodic grid at all, so the derivative operators must be rebuilt from
boundary-aware differences rather than taken off the shelf. The way out is
structural rather than empirical. Because the network is trained to return the
state at $t_k+\Delta t$ from the state at $t_k$, $\partial_t V_m$ in
\eqref{eq:ch35-residual} can be replaced by a difference of successive predicted
frames, so that no spectral derivative of a front is ever formed; the spatial
term $\nabla\cdot(\mathbf{D}\nabla V_m)$ still requires the boundary-aware
differences just described. The objection above is therefore aimed at the naive
reading of \eqref{eq:ch35-loss}, in which the residual is taken spectrally
through the upstroke, and not at the formulation this chapter uses.

\index{zero-shot super-resolution!in a cardiac surrogate}The same observation
undercuts the cleanest-sounding claim one could make for a surrogate in this
domain, and it is worth stating the gap exactly. Proposition~\ref{prop:ch06-zeroshot} upgrades
resolution transfer from an empirical observation to an identity, but its
hypothesis is that the input is band-limited within the retained modes $Z_{k_{\max}}$,
and its conclusion is that the \emph{spectral branch} of a truncated FNO
reproduces the same trigonometric polynomial on a finer point set. Cardiac
excitation violates the hypothesis at the front and satisfies it in the
quiescent tissue between waves. So the report is a divided one, and this
book can take it no further than that: the transfer is licensed by
\eqref{eq:ch06-zeroshot} in the quiescent tissue, and across the front it is an
assumption for which no measurement is cited here. Two conditions any such
measurement would have to satisfy belong in the statement in advance, because
they are the two that a favorable number most easily conceals. The fine grid must
lie beyond the grid the ground-truth data were produced on, since evaluating on
the data grid tests interpolation of the labels rather than transfer; and the
sweep must be scored inside the recursive roll-out an autonomous simulator would
have to run, not point-to-point with the true input supplied at every step. A
factor of $M$ in linear resolution on a two-dimensional tissue
preparation is a factor of $M^2$ in stored nodes, which is a different
statement from the one \eqref{eq:ch06-zeroshot} licenses and a more useful one
clinically, since lesion-size estimates are area-type quantities.

Panel~(b) of
Figure~\ref{fig:vol3_sub}, printed in
\chapref{ch:part3-subsurface-multiphase-flow-and-carbon-stor}, shows the
geometry this map is eventually being asked to reproduce: a reentrant spiral on
a ventricle-shaped body, colored by $V_m$. That plate pairs it with the
subsurface problem of the previous chapter because the two share one obstruction
--- a sharp front on a body the spectral basis was not written for. No cardiac
operator result on an anatomical mesh is cited anywhere in this book; the panel
therefore depicts the problem this chapter formulates, not a result anyone has
reported.
\index{rotor (reentrant wave)}The reader should take two things from panel~(b).
The first is topological: a spiral is
not a traveling wave bent into a circle but a solution with a core where $V_m$ is
undefined in phase, tissue that is depolarizing while its neighbor is still
refractory, and the rotation period at the core is typically faster than the
surrounding cycle length. Reentry of this kind self-sustains: nothing needs to
trigger it, and only an intervention---a second wave colliding with the front, a
defibrillating shock that recruits most of the tissue simultaneously---ends it.
The second is about the surrogate. Everything in the picture that makes reentry
hard to simulate, the thin front, the sharp gradient at the core, the long
duration, is the same thing that makes \eqref{eq:ch35-residual} hard to evaluate
and \eqref{eq:ch06-zeroshot} inapplicable. Panel~(b) is thus the visual form of
the objection in the previous paragraph, which is the reason the two belong
together in one section.
\index{reentry!spiral wave}\index{phase singularity}

% -----------------------------------------------------------------------------
\section{Documented Claims and Runtime Accounting}
\label{sec:ch35-claims}
\index{arrhythmia!surrogate verification}
% -----------------------------------------------------------------------------

Three findings reach this chapter attached to the cardiac PINO --- roll-out
fidelity over extended horizons, zero-shot resolution transfer, and a reduction
in simulation time --- and separating what a primary study supports from what the
inherited framing asserts was the task this section was written for. The separation comes
out all one way: this book cites no study that reports a physics-informed operator
surrogate for cardiac tissue, so the three findings stand as assertion carried
forward rather than as cited results. What survives is the formulation of
Section~\ref{sec:ch35-map}, the classical cost side, and the arithmetic by which
any future measurement would be judged. Each finding is worth walking through in
that light, because each has a precise form the book can state even where it
cannot state a number.

Take roll-out fidelity first. In its strongest legitimate form the claim is that
recursive rollouts---predictions fed back as
inputs---hold their accuracy over a simulated horizon, and that a model shown only
planar and centrifugal waves reproduces the stable spiral. That is a statement
about a learned dynamical system rather than about a regression fit. The rollout
arithmetic of\index{autoregressive rollout!cardiac horizons}
\eqref{eq:ch09-onesided} prices it: with one-step error $\tau$ at the true
states and Lipschitz constant $\Lambda_\theta$, pathwise error obeys
\eqref{eq:ch09-gronwall}, so sustained accuracy over $K$ steps is evidence that
the trained map's local expansion along the realized trajectory is not large on
the tested horizon---and nothing more. No term in \eqref{eq:ch35-loss} bounds
$\Lambda_\theta$; \eqref{eq:ch09-lipschitz} makes it an explicit functional of
the weights, and \chapref{ch:part1-open-theoretical-problems-and-mathematical}
argues that requiring $\Lambda_\theta<1$ would be structurally wrong for
dynamics whose leading Lyapunov exponents are positive. Excitable tissue is a
case in point: a surrogate that contracted every distance could not sustain a
rotor. Long-horizon stability is therefore an observed property of specific
trained models, never a consequence of the objective, and here it is an
observation this book has no data for.

The taxonomy such an observation would be scored against is worth keeping
regardless, because it is a taxonomy of the dynamics rather than of any study.
Four regimes organize monodomain wave simulations: planar fronts, centrifugal
spread from a focal stimulus, a stable reentrant spiral, and spiral break-up into
several filaments, the last two initiated by cross-stimulation rather than by
pacing. The first two are conduction; the third is the sustained arrhythmia read
clinically as monomorphic tachycardia; the fourth is the polymorphic and
fibrillatory case this chapter opens with. Break-up is where a soft residual
penalty is least likely to suffice, for the reason the previous paragraph
supplies: wavelet collision, core meander and the nucleation of new fronts are
events at the resolution limit, and \eqref{eq:ch09-gronwall} bounds error growth
only along a trajectory on which the map is already known to be mildly expanding.
Roll-out fidelity for a stable spiral and roll-out fidelity through break-up are
therefore different claims, and an average relative $L^2$ error in $V_m$ is weaker
evidence for the second than for the first. This book records the distinction as
a requirement on any future measurement, not as a result.

The second item is resolution transfer, with the conditions it would have to meet
stated in Section~\ref{sec:ch35-machinery}; no measurement of it in this domain is cited
here. The third is cost. The inherited figures present a cardiac PINO as a step toward
efficient and scalable simulation, motivated by planning catheter
ablation\index{catheter ablation} and
device implantation; the motivation is legitimate and the workflow it points at is
not validated anywhere in this book's record. What accompanies that framing is a
pairing---minutes of finite-element solve to sub-second inference---and a
multiplier of about $1{,}500\times$ set against it. Neither side of that pairing
is a measurement this book can attribute. The solve side is the
order-of-magnitude figure of Section~\ref{sec:ch35-monodomain}, carried without a primary
reference; the inference side has no reference behind it at all.
\chapref{ch:part3-synthesis-matrix-and-real-world-deployment}
reaches the same verdict in its cardiac row, printing no multiplier where no pair
of timings exists. Table~\ref{tab:ch35cost} lays the accounting out
row by row.
\begin{table}[t]
\centering
\small
\caption{Cost and evidence accounting for the monodomain problem and for the
physics-informed operator formulation \eqref{eq:ch35-loss} of it, built on the
PINO template \cite{li2021physics}. ``inherited'' marks values carried forward by
this Part with no primary reference behind them; ``this book'' marks
derivations made here; ``not documented'' marks quantities this book cannot
attribute. The conclusion to draw is that the classical column is
order-of-magnitude bookkeeping, the surrogate column is a formulation with no
measurements attached to it in this record, and the two rows that would decide
clinical use---stability through break-up, and any runtime gain---are the empty
ones.}
\label{tab:ch35cost}
\begin{tabularx}{\textwidth}{>{\raggedright\arraybackslash}p{2.7cm}>{\raggedright\arraybackslash}X>{\raggedright\arraybackslash}X>{\raggedright\arraybackslash}p{1.75cm}}
\toprule
\textbf{Quantity} & \textbf{Classical finite-element solve} & \textbf{PINO surrogate} & \textbf{Basis} \\
\midrule
Spatial mesh & $h\sim 200\ \mu\text{m}$, tetrahedral, body-fitted & grid not fixed by the architecture; no evaluation resolution measured in anything cited here & inherited / not documented \\
Time step & $\Delta t\sim 10\ \mu\text{s}$, small and adaptive, accuracy-limited by the upstroke and \eqref{eq:ch35-cfl} & one application of $\mathcal{G}_{\Delta t}$ per step; frame spacing a design choice, not a measurement & inherited / not documented \\
Supervision & none required: the model is integrated directly & reference trajectories plus the cell model \eqref{eq:ch35-ion}--\eqref{eq:ch35-recovery} as a soft residual \eqref{eq:ch35-residual}, with a boundary term for \eqref{eq:ch35-no-flux} & this book \\
Refinement cost & $\mathcal{O}(h^{-(d+2)})$ by \eqref{eq:ch35-cost}: re-mesh, re-solve, re-verify & re-evaluate on a finer point set, no retraining; licensed by \eqref{eq:ch06-zeroshot} only for band-limited inputs & this book \\
Horizon and stability & stable under \eqref{eq:ch35-cfl} and solver controls & governed by $\Lambda_\theta$ through \eqref{eq:ch09-gronwall}; no bound in \eqref{eq:ch35-loss}, no roll-out measurement cited here, and break-up the least likely case to hold & not documented \\
Runtime gain & minutes-scale per solve in the framing this Part carries, with no primary reference & no pair of timings on the record; the $\approx1{,}500\times$ factor is secondhand arithmetic, not a measurement & inherited / not documented \\
Clinical qualification & verified against the discretized model & no prospective, multi-center patient validation & not documented \\
\bottomrule
\end{tabularx}
\end{table}
Read as a whole, the table's message has changed register from the one the
findings it tabulates were written to support. The classical column is order-of-magnitude bookkeeping that
carries no primary reference; the surrogate column is a formulation rather than a set of
measurements; and the two rows that would decide
clinical use, stability through break-up and any runtime gain, are the empty ones.
That is not an argument that a physics-informed operator cannot do this problem.
It is a statement of what has to be measured before the phrase ``cardiac
surrogate'' carries content.

% -----------------------------------------------------------------------------
\section{Three Clinical Constraints on Any Cardiac Answer}
\label{sec:ch35-clinical}
% -----------------------------------------------------------------------------

\index{model verification!versus physiological validation}
\index{out-of-distribution (OOD) failure!clinical surrogate}
Three domain constraints decide how far any answer to that question travels, and
none of them is a network-architecture matter. The first is the target of
verification. A surrogate trained on solver output is accurate to the solver,
which means it inherits the solver's model: the reaction kinetics
\eqref{eq:ch35-ion}, the diffusion tensor \eqref{eq:ch35-dtensor} --- which
collapses to a scalar diffusivity in the two-dimensional setting this chapter can
write down completely, and which is exactly the setting a verification would be
carried out in --- the zero-flux boundary
\eqref{eq:ch35-no-flux}, and the assumption that the monodomain reduction
suffices. Agreement with a finite-element solver of the openCARP
type---the class of simulator that would generate such training data---establishes
that the operator has learned a discretized two-variable model. It
says nothing about whether that model
describes the patient, and the parameters that would matter clinically---the
longitudinal-to-transverse conductivity ratio, the restitution slope, the
location of interstitial fibrosis---are among the least measurable quantities in
the organ. Diffusion-tensor magnetic resonance imaging
(DT-MRI)\index{diffusion-tensor magnetic resonance imaging (DT-MRI)} supplies $\mathbf{f}_0$
on image voxels coarser than the mesh the equation wants, and
transmural fiber rotation and laminar sheet structure fall below the resolution
of the acquisition; the model's anisotropy is therefore an interpolation of an
interpolation.

The second constraint is the training measure. \chapref{ch:part1-operator-approximation-theory-and-universa}
bounds approximation error on the compact support of the measure $\mu$ from
which inputs are drawn, and offers nothing outside it. In a cardiac study $\mu$
is a family of stimulus patterns on a family of geometries, generated by a
simulator on a grid of chosen resolution with a chosen parameter range, and in
this chapter's formulation that range is whatever the designer declares, since no
cardiac training measure is supplied anywhere in this book. The
experiment that would settle the question is a held-out one: train on
conductivities in a narrow band around the baseline and score the surrogate
outside the band. That is the only direct measurement of what lies off the
support, and this book cites no such measurement for cardiac tissue, so the
penalty for leaving the band is unknown rather than small. A
ventricle with infarct scar, a stimulus delivered during the vulnerable period,
or a rotation rate outside the training distribution is an input off that
support, and the field that comes back will look like a wave regardless. This
is why the warning recorded in
\chapref{ch:part2-production-deployment-and-high-throughput-}---that an
out-of-distribution field returns a plausible field rather than an error
message---is not an optional accessory in this domain: a surrogate that cannot
refuse is worse than a slow one, and
\chapref{ch:part3-synthesis-matrix-and-real-world-deployment} is where this book
takes up the uncertainty machinery that would have to supply the refusal.

The third is the pathway. Anything used to guide an ablation lesion set is
software as a medical device, and the requirements this book
collects under regulatory clearance---prospective, multi-center evaluation on
diverse cohorts; documented verification and validation; traceability of the
model to the measured inputs---are described in
\chapref{ch:part3-synthesis-matrix-and-real-world-deployment}. What this chapter
adds is the specific reason the barrier is higher here than in aerodynamics: the
quantity of interest is not a field average but the location and lifetime of a
phase singularity, an event-type output for which relative $L^2$ error in $V_m$
is a weak proxy. The proxy is weaker still when the error is scored on the voltage
channel alone and the recovery variable $w$ is left latent, which is the natural
consequence of an output layer read as a depolarization map rather than as the
state pair of \eqref{eq:ch35-timemap}; and the training
labels are simulations of a model whose
parameters are the very things a clinical study would want to infer.

\takeaway{What the Cardiac Surrogate Is Entitled to Claim}{A physics-informed
FNO applied to \eqref{eq:monodomain} with Aliev-Panfilov kinetics as a soft
residual \eqref{eq:ch35-loss} is a well-posed proposal. The reaction term is
known exactly even where the tissue parameters are not, the residual
\eqref{eq:ch35-residual} is evaluable at any resolution the storage supports, and
\eqref{eq:ch09-gronwall} states precisely what a sustained roll-out would and
would not demonstrate. What the formulation is not entitled to, on this book's
record, is a single measured claim: not roll-out fidelity through break-up, not
resolution transfer across a front, not anisotropy or anatomy carried as an
argument of the operator, not a speedup factor, not a bound on the rollout
Lipschitz constant, not physiological validation beyond the discretized model it
would be trained on. In Table~\ref{tab:ch35cost} the classical column is
textbook bookkeeping and the surrogate column awaits its first measurement; the
accounting is the contribution, because it is the part that survives whichever
study fills the rows in.}

% -----------------------------------------------------------------------------

The cardiac case and the tomographic case are usually presented as separate
application areas, and the accounting just given explains why they belong side by
side. Both put a neural operator in contact with a continuous forward map that
the practitioner cannot invert exactly; both depend on resolution-agnostic
evaluation, since the question a radiologist or a cardiologist asks is about a
field at a scale the training grid never carried; and both would be trained on
targets produced by an algorithm rather than measured---a finite-element solve
there, a reconstruction convention here. The direction of the error differs,
though, and it is the difference that structures the next chapter. In the
monodomain problem the operator is well posed and the difficulty is cost; in
sparse-view \emph{Computed Tomography} (CT) the forward problem itself is the
obstacle, and the classical \emph{Filtered Backprojection} (FBP) inverse is
replaced by an ill-posed sampling of the Radon transform whose artifacts---streaks,
blurred boundaries---are not solver error but missing information.
\chapref{ch:part3-medical-imaging-and-tomographic-reconstruc} develops that
failure mode with the numerical detail the operator treatment requires, and by
then the accounting discipline of Table~\ref{tab:ch35cost} will have a second
data set to be applied to.

% -----------------------------------------------------------------------------
